\documentclass[11pt,letterpaper]{article}
\usepackage[margin=1in]{geometry}
\usepackage{fontspec}
\usepackage{xeCJK}
\setmainfont{Times New Roman}
\setCJKmainfont[AutoFakeBold=true]{DFKai-SB}
\usepackage{setspace}
\usepackage{enumitem}
\usepackage[para]{footmisc}
\renewcommand{\footnotesize}{\fontsize{8}{9.2}\selectfont}
\setlength{\skip\footins}{0.5em plus 0.1em minus 0.1em}
\usepackage{hyperref}
\hypersetup{colorlinks=true, allcolors=black}
\setlength{\parindent}{0pt}
\setlength{\parskip}{0.38em}
\setstretch{1.0}
\pagestyle{empty}

\begin{document}

\begin{center}
{\Large 研究興趣與未來研究計畫}\\[0.25em]
許可證（Ke Cheng Hsu）\\[0.15em]
{\small 應徵中華經濟研究院助研究員}
\end{center}

\noindent\textbf{研究動機。}線上市場的運作，依賴一套關於「反悔」的規則：消費者可以退回買下的商品，書店可以退回賣不掉的存貨，平台可以在交易未完成時退還費用。誰制定這些規則、規則該如何設計，決定哪些商品被買走、哪些資源被浪費，以及成本由誰承擔。我以實證方法研究這些問題。

\medskip
\noindent\textbf{研究興趣。}我的研究領域為實證產業組織與數位經濟。我研究的，是在不確定性尚未消除前就必須做出、而事後調整成本高昂的決策。消費者不確定產品價值，可以在付款前透過搜尋了解，也可以在付款後透過退貨了解；賣方不確定需求，必須預先定價與備貨。一項看似修正某一邊際的政策，往往扭曲另一邊際，而這兩個邊際是同時決定的。這條「聯合決定」的主線，貫穿我對產品退貨、消費者搜尋，以及需求不確定性下庫存決策的研究，也是我帶入政策分析的視角。

\medskip
\noindent\textbf{工作論文。}
\begin{itemize}[leftmargin=1.6em,itemsep=0.2em,topsep=0.1em,parsep=0.1em]
\item \textbf{Designing Return Policy in E-Commerce: A Trade-off between Purchase and Return Efficiency}（求職論文）。退貨政策決定兩個邊際：消費者是否購買，以及收到商品後是否退貨。較嚴格的政策能抑制購買後無效率的退貨；但退貨減少後，賣方在每筆成交中保留更多，較高的價格變得更有利可圖，可能把原本有價值的購買擠出市場。我使用淘寶一家女裝零售商的交易資料，估計一個涵蓋六個服裝類別、獨占定價下的購買、退貨保險與退貨決策結構模型。結果顯示：利潤率與殘值較高時，補貼退貨運費的寬鬆政策較能擴大有效率的購買；利潤率較薄時，較嚴格的政策較能抑制過度退貨；而折扣價搭配適度退貨手續費的選單，能平衡購買與退貨效率，在多數類別中優於前兩者。台灣的法規讓這個問題格外具體：《消費者保護法》第19條賦予網購消費者七日猶豫期，可無理由免費退貨。\footnote{消費者保護法第19條。英文條文：全國法規資料庫，\url{https://law.moj.gov.tw/ENG/LawClass/LawAll.aspx?pcode=J0170001}。}免費退貨既是法定預設，賣方在這項制度下應如何定價、如何設計退貨機制，便是受管制市場中的現實問題，而非紙上假設。
\item \textbf{Product Design with Heterogeneous Search Costs}（與Pak Hung Au合著）。我們在序貫搜尋模型中引入搜尋成本異質性。低搜尋成本的消費者尋找利基產品，高搜尋成本的消費者購買標準化產品。當搜尋成本下降，更多消費者轉向搜尋利基產品，這種「誰在搜尋」的組成改變，可能同時推高兩類產品的價格。
\end{itemize}

\medskip
\noindent\textbf{進行中的研究。}
\begin{itemize}[leftmargin=1.6em,itemsep=0.2em,topsep=0.1em,parsep=0.1em]
\item \textbf{Pass-Through and Incidence with Product Returns}（與Ce Bian、Liang Zhao合著）。本計畫探討：當購買並非最終交易時，租稅如何轉嫁。只有當買方保留商品，賣方才真正收到價款，因此每筆銷售都承擔生產成本、退貨處理成本，以及退貨時無法回收的租稅或費用。價格也會影響退貨率，退貨率再回饋到利潤率與供給量，因此轉嫁可能更強、更弱，甚至為負。忽略退貨邊際的計算，可能逆轉預測的價格與福利效果。本計畫也探討保留的銷售與退貨應如何課稅：實務上營業稅在退貨時全額退還，平台轉介費卻只部分退還。台灣的電商平台讓這個問題格外現實：蝦皮分階段調漲賣家費用，合計抽成據報導已超過14\%，並於2026年調高金流與系統處理費、要求所有賣家加入免運專案；同時，消費者在猶豫期內可免費退貨，平台對退貨訂單按比例退還成交手續費。\footnote{蝦皮購物賣家幫助中心：成交手續費與金流費，\url{https://seller.shopee.tw/edu/article/26470}；2026年1月費率更新與免運專案，\url{https://seller.shopee.tw/edu/article/26430}；退貨退款訂單不收成交手續費、費用按比例退還，\url{https://seller.shopee.tw/edu/article/12157}。合計抽成超過14\%與賣家反彈：《天下雜誌》專訪蝦皮台灣總經理，2026年，\url{https://www.cw.com.tw/article/5141245}。}這些費用如何轉嫁到價格、各項費用在退貨時的退還方式如何形塑轉嫁，是台灣平台與賣家之間持續中的爭議。
\item \textbf{Too Much Unsold Stock?: Environmental Consequences of Price Commitment}（與Kohei Kawaguchi、Kei Kawai合著）。在日本圖書市場，出版社在維持轉售價格制度下訂定零售價，書店將未售出的書籍退回打成紙漿。固定零售價鼓勵書店在需求轉弱時仍持有庫存，這有助於競爭；但在需求低迷時也無法降價出清，滯銷存貨於是堆積，造成環境成本。我們建立含需求不確定性的動態寡占模型，以書店月度資料估計，探討有多少滯銷存貨源自價格承諾，以及押金退還稅能如何減少浪費。台灣《公平交易法》將維持轉售價格視為原則違法，公平交易委員會曾處分固定書籍轉售價格的出版商，包括對天下遠見出版將《賈伯斯傳》售價固定為新臺幣599元裁罰新臺幣20萬元，以及對龍騰限制書店參考書折扣裁罰新臺幣50萬元。\footnote{公平交易委員會處分：公處字第101137號（2012年9月19日，天下遠見出版）；（89）公處字第207號（2000年11月30日，龍騰）。案例摘要：\url{https://www.apeccp.org.tw/htdocs/doc/Taipei/Case/dec2012/1136-1089case07.html}；\url{https://www.apeccp.org.tw/htdocs/doc/Taipei/Case/D0505200.htm}。}
\end{itemize}

\medskip
\noindent\textbf{未來研究計畫。}未來五年，我將把這個研究議程發展為對「造成浪費的市場規則」更廣泛的分析，並連結環境與循環經濟政策。產品退貨與滯銷存貨，是同一個市場失靈在銷售前後兩個階段的表現，兩者都受價格、費用與交易反悔權的形塑。我看到三個方向。

第一，我將一併研究搜尋、購買與退貨。搜尋與退貨是解決同一個媒合不確定性的兩種方式，一個在付款前、一個在付款後，因此互為替代。更好的搜尋會降低買方帶入購買的殘餘不確定性，理應降低退貨率；但它也讓賣方可以收取更高價格，因為買方對外部選項有更充分的資訊。哪股力量占主導，決定平台的推薦系統是真的降低退貨，還是只是推高價格。我希望估計這個聯合決策，並從中還原搜尋強度與退貨率之間的替代關係。可能的成果是一套福利核算，將媒合品質的提升與價格上升分開，並檢驗旨在提升媒合品質的推薦系統，是否反而把錯配的成本轉嫁到退貨運輸與處置上。在台灣，免費退貨是法定預設，\footnote{消費者保護法第19條，\url{https://law.moj.gov.tw/ENG/LawClass/LawAll.aspx?pcode=J0170001}。}搜尋與退貨之間的替代、以及平台推薦在其中的角色，正是消費者已經面對的市場特徵。

第二，我將把退貨與滯銷存貨作為環境問題來研究，並評估減少浪費的政策工具。在服裝業，消費者一次訂購數個尺寸或顏色，留下一件、退回其餘，退回的商品價值大幅折損；在出版業，固定零售價產生被打成紙漿的滯銷書籍。兩者都是保護某一邊際的規則，在另一邊際製造實體浪費，因為制定規則的一方不承擔處置成本。我計畫量化這些浪費，並比較能將其內部化的工具，包括退貨手續費、退貨運費、押金退還稅與價格彈性規則。直觀上，每種工具都把退貨的私人成本推向社會成本，但它們在多大程度上同時抑制有價值的購買並不相同，因此這項比較將排定：每犧牲一單位交易，哪種工具減少的浪費最多。台灣環境部已管制電商包裝：禁用PVC、規定紙箱的再生紙含量達90\%、塑膠包材的再生料含量達25\%，並推動大型業者達成包裝減重與重複使用目標；退貨則在這些流量之上，再增加逆向物流與包裝廢棄物。\footnote{環境部《網際網路購物包裝限制使用對象及實施方式》，2023年7月1日施行。英文摘要：\url{https://www.moenv.gov.tw/en/news/press-releases/15507.html}。}隨著台灣延伸生產者責任與循環經濟規範趨嚴，這項研究可為電商退貨與滯銷存貨應如何計量、收費或抑制提供依據。

第三，我將研究退貨與退款規則如何形塑平台的雙邊市場。退貨政策是影響參與的槓桿，不只影響加入後的行為。寬鬆的政策吸引更多買方，包括較可能退貨的買方；這提高賣方面對的退貨率、降低賣方收益，於是部分賣方退出，存活者重新定價。跨邊效果因此可能抵銷寬鬆政策原本要創造的部分需求增益。我計畫估計一邊的組成變化如何回饋到另一邊，並比較平台會選擇的退貨政策與監理機關會選擇的退貨政策。這將建立從跨邊網路效應、到退貨運費與退貨手續費在買方、賣方與平台間最適分攤的對應關係。公平交易委員會的數位平台競爭研究，已使平台為其雙邊制定的規則在台灣受到檢視。\footnote{公平交易委員會委託研究，中華經濟研究院執行：《數位平臺產業現況、發展趨勢及競爭規範—以我國為中心》（民國111年），\url{https://www.ftc.gov.tw/upload/1f59ab5d-8cae-410a-a8b1-2080fdab977c.pdf}。}

\medskip
\noindent\textbf{對中華經濟研究院的貢獻。}我的領域是實證產業組織，而我的工具正是貴院已經在使用的工具。中華經濟研究院接受公平交易委員會委託進行競爭研究，包括網路銷售市場的實證評估，以及數位平台與競爭規範的研究。\footnote{公平交易委員會委託研究，均由中華經濟研究院執行：《網路銷售市場競爭評估之實證研究》（民國109年），\url{https://www.ftc.gov.tw/upload/885e4488-950e-4cf9-a774-6a224b8df58c.pdf}；《數位平臺產業現況、發展趨勢及競爭規範—以我國為中心》（民國111年），\url{https://www.ftc.gov.tw/upload/1f59ab5d-8cae-410a-a8b1-2080fdab977c.pdf}。}這些報告界定相關市場、評估多邊網路效應與市場力量，並檢視獨家交易、維持轉售價格、最惠客戶條款與平台平價等行為。這正是我研究的領域。我的訓練是處理豐富的微觀資料，估計需求、供給與廠商行為的結構模型，並用這些模型評估規則或價格的改變如何影響市場力量、消費者剩餘與福利。我將把這個能力帶進貴院的競爭與數位經濟研究。

我的數個計畫與上述工作直接相關：轉嫁計畫探討在購買可以退貨時，費用與租稅如何傳導到價格；滯銷存貨計畫估計維持轉售價格的競爭效果，這正是公平交易委員會點名的垂直交易限制之一；求職論文與搜尋計畫則處理數位平台研究關注的定價、搜尋與平台設計問題。在這些工作中，我能完成從原始交易或掃描資料、到可供部會採納的政策模擬的完整流程。

同一套工具也服務於貴院其他中心。對台灣產業、能源及綠色經濟相關中心而言，我在滯銷存貨與產品退貨上的研究，衡量定價規則與運費補貼造成的浪費，並評估可減少浪費的工具，如押金退款稅與補貨費。對經濟預測而言，我的轉嫁研究顯示，忽略退貨邊際可能逆轉對成本衝擊的預測價格反應。我在台灣長大，能以中、英文處理國內資料，並與利害關係人合作。我在中華經濟研究院的目標，是把結構估計轉化為競爭、數位平台、租稅與環境政策上清楚、有實證依據的建議。

\end{document}
